% AI-written, AI cross-reviewed; no human sign-off; not a claim that RH is proved
% Atlas unit P040. Local manuscript v3, item 7.17, heading on printed PDF p. 82.
% UpperTheory.V.LDataCardLeFinrank.vFormula_of_classicalTheorems
% paper2/inputs_one.tex:730–771. Verbatim excerpt; SHA-256 87c190ce70fdd9a242449049f7f9454c752469950e8cdcd41edd2658fda79674.
% Original TeX uses the macros and cross-reference keys of the source manuscript.


% Start of atlas edition annotation. The following block is outside the original excerpt.
\paragraph{Status of named inputs in this reading.}
\noindent\nolinkurl{GlobalTheory.Published.PoitouTateDuality} --- hypothesis at this head (\path{Development/AbelianVarietyTheory/PublishedTheorems/GaloisCohomology₂/ExistsPrimitiveRootPow.lean:13--27}); not compared.\par
\noindent\nolinkurl{GlobalTheory.TateLimit.local_tate_duality_h2_limit} --- in-repo constant (\path{Development/AbelianVarietyTheory/Proofs/FoundationProofs/LengthEqLengthInvariantsOfTorsionLevels.lean:295--302}); not compared.\par
% End of atlas edition annotation.

\begin{formalstatement}[A compact-coefficient length formula]\label{pn:P040}
Let \(F\) be a number field and \(\ell\) a prime such that \(\ell\ne2\)
or \(F\) has no real place. Assume finite-coefficient Poitou--Tate
duality, local finiteness and local Tate duality. For every allowed
set of local data \(D\) and every \(e\in\mathbb Z\), with
\(A_e=(\mathbb Q_\ell/\mathbb Z_\ell)(-e)\),
\[
 \operatorname{length}H^2(G_{F,S},\mathbb Z_\ell(e+1))
      +\nu_{\rm Gal}(\ell,F,|e|)
   =\operatorname{length}\Sha^1(D,A_e)
      +\sum_{w\in D.S}\operatorname{length}H^0(F_w,A_e).
\]
All lengths are extended \(\mathbb Z_\ell\)-lengths;
\(\Sha^1\) is the kernel of the specified localization.
\end{formalstatement}\nodestatus{P040}
\begin{outline}
At level \(j\ge1\) take \(T_j=(\mathbb Z/\ell^j)(e+1)\) and its
Cartier dual \(B_j=A_e[\ell^j]\). Write \(K_j\) for the localization
kernel on \(H^1(G_{F,S},B_j)\) and a superscript \(\vee\) for the
character dual. The last maps of Poitou--Tate and local duality yield
\[
 0\longrightarrow K_j^\vee\longrightarrow H^2(G_{F,S},T_j)
 \longrightarrow\bigoplus_{w\in D.S} H^0(F_w,B_j)^\vee
 \longrightarrow H^0(G_{F,S},B_j)^\vee\longrightarrow0 .
\]
The modified archimedean terms vanish in the indicated case. All
groups at finite level are finite, so images stabilize and this
sequence remains exact under inverse limits. On the discrete side
a continuous cochain has finite image, and filtered union commutes
with the cohomology and localization kernels in question. On the
compact side the cochain reductions are surjective (lift on each
clopen fibre); the possible derived limit of finite \(H^1\)'s
vanishes. Thus the middle term becomes compact \(H^2\), and the
other terms are the duals of the indicated discrete groups.
Additivity of extended length, its preservation under this duality,
and the identification of the global invariant length with
\(\nu_{\rm Gal}(\ell,F,|e|)\) give the formula. This is the
finite-to-compact argument of Blueprint Z2.K8~\bploc{lem:P-r3-k8}. For the finite
duality statements see J.~S.~Milne, \emph{Arithmetic Duality Theorems},
I.2, Corollary~2.3 and I.4, Theorem~4.10.
\cite{r21-adt}
\end{outline}
